%% ============================================================
%%  SECCIÓN 4 — Fundamento Teórico
%%  Responsable: [Nombre]
%% ============================================================

\section{Fundamento Teórico}

\subsection{Elasticidad y deformación}

Cuando un cuerpo sólido es sometido a una fuerza externa, experimenta un cambio en su
forma o dimensiones denominado \textit{deformación}. Si el cuerpo recupera su geometría
original al retirar la fuerza, se dice que el comportamiento es elástico. La propiedad
que caracteriza esta capacidad de recuperación se denomina \textit{elasticidad}
\cite{serway}.

La deformación longitudinal unitaria $\varepsilon$ cuantifica el cambio relativo en la
longitud del cuerpo respecto a su longitud natural $\ell_0$:
\begin{equation}
    \varepsilon = \frac{\Delta\ell}{\ell_0} = \frac{\ell - \ell_0}{\ell_0},
    \label{eq:deformacion_unitaria}
\end{equation}
donde $\ell$ es la longitud bajo la acción de la fuerza y $\Delta\ell$ es la deformación
total. Al ser un cociente de longitudes, $\varepsilon$ es adimensional.

\subsection{Esfuerzo de tensión}

El esfuerzo $\sigma$ representa la fuerza interna por unidad de área que actúa sobre la
sección transversal del material. Se distinguen dos definiciones según el área de
referencia empleada \cite{halliday}:

El \textit{esfuerzo técnico} (o de ingeniería) utiliza el área de la sección transversal
original $S_0$, medida antes de aplicar la carga:
\begin{equation}
    \sigma_0 = \frac{F}{S_0} \quad (\unit{\pascal}),
    \label{eq:esfuerzo_tecnico}
\end{equation}

El \textit{esfuerzo real} utiliza el área deformada $S$ de la sección transversal bajo
la acción de la fuerza $F$:
\begin{equation}
    \sigma = \frac{F}{S} \quad (\unit{\pascal}).
    \label{eq:esfuerzo_real}
\end{equation}

Para deformaciones pequeñas, $S \approx S_0$ y ambas definiciones son equivalentes.

\subsection{Ley de Hooke y módulo de Young}

En 1658, Robert Hooke estableció empíricamente que, para deformaciones dentro del límite
elástico, el esfuerzo es proporcional a la deformación unitaria. Esta relación se conoce
como la \textit{Ley de Hooke} y se expresa como \cite{serway}:
\begin{equation}
    \sigma = Y\varepsilon,
    \label{eq:hooke}
\end{equation}
donde la constante de proporcionalidad $Y$ es el \textit{módulo de Young} o módulo de
elasticidad en tracción, con unidades de \unit{\pascal}. El módulo de Young caracteriza
la rigidez del material: cuanto mayor es $Y$, mayor es la resistencia del material a
deformarse.

Combinando las Ecs.~\eqref{eq:deformacion_unitaria}, \eqref{eq:esfuerzo_tecnico} y
\eqref{eq:hooke}, se obtiene para una barra prismática sometida a tracción axial uniforme:
\begin{equation}
    F = \frac{Y S_0}{\ell_0}\,\Delta\ell = k\,\Delta\ell,
    \label{eq:hooke_resorte}
\end{equation}
donde $k = Y S_0 / \ell_0$ es la \textit{constante recuperadora}
(\unit{\newton\per\metre}), que cuantifica la rigidez del elemento elástico.

Es importante señalar que esta deducción es estrictamente válida para barras rectas de
sección uniforme bajo deformación axial homogénea. En el caso de un resorte helicoidal,
la deformación longitudinal observada surge predominantemente de la \textit{torsión del
alambre} que lo compone, no de tracción axial directa. Por ello, la expresión
$F = k\,\Delta\ell$ se aplica al resorte como una analogía macroscópica formal: $k$
sigue siendo la constante recuperadora del sistema, pero no puede relacionarse
directamente con $Y$ mediante $k = Y S_0/\ell_0$. El modelo riguroso del resorte
helicoidal proviene de la teoría de torsión \cite{halliday}:
\begin{equation}
    k = \frac{G\,d^4}{8\,D^3\,N},
    \label{eq:k_resorte_helico}
\end{equation}
donde $G$ es el módulo de rigidez del material, $d$ el diámetro del alambre, $D$ el
diámetro medio de la hélice y $N$ el número de espiras activas.

\subsection{Límite elástico y diagrama esfuerzo--deformación}

La Ley de Hooke es válida únicamente para deformaciones pequeñas, dentro de un intervalo
denominado \textit{región elástica}. Superado el \textit{límite elástico} $\sigma_{\text{límite}}$,
la relación lineal deja de cumplirse y la deformación se vuelve permanente
\cite{halliday}.

En la Fig.~\ref{fig:diagrama_sigma_epsilon} se muestra el diagrama $\sigma$--$\varepsilon$
típico de un material dúctil, donde se identifican la región elástica lineal, el límite
elástico y el punto de ruptura.

\begin{figure}[H]
    \centering
    \includegraphics[width=0.55\linewidth]{fig_diagrama_hooke}
    \caption{Diagrama esfuerzo vs.\ deformación unitaria de un material dúctil.
             Se distinguen la región elástica lineal (Ley de Hooke), el límite elástico
             $\sigma_{\text{límite}}$ y el punto de ruptura.}
    \label{fig:diagrama_sigma_epsilon}
\end{figure}

\subsection{Histéresis elástica}

En materiales no perfectamente elásticos, como la liga o el jebe, las curvas de carga y
descarga en el diagrama $\sigma$--$\varepsilon$ no coinciden, formando un lazo cerrado
denominado \textit{lazo de histéresis}. El área encerrada por este lazo representa la
energía disipada por el material en un ciclo completo de deformación, asociada a
mecanismos de disipación interna como la viscosidad y la fricción molecular
\cite{serway}.

\subsection{Trabajo de deformación}

El trabajo realizado por la fuerza $F$ para deformar el resorte desde su longitud natural
hasta una elongación $\Delta\ell_f$ se obtiene integrando la Ec.~\eqref{eq:hooke_resorte}:
\begin{equation}
    W = \int_0^{\Delta\ell_f} F\,\mathrm{d}(\Delta\ell)
      = \int_0^{\Delta\ell_f} k\,\Delta\ell\,\mathrm{d}(\Delta\ell)
      = \frac{1}{2}k\,(\Delta\ell_f)^2.
    \label{eq:trabajo_resorte}
\end{equation}
Experimentalmente, este trabajo se puede estimar como el área bajo la curva $F$
vs.\ $\Delta\ell$ mediante integración numérica (regla del trapecio).

\subsection{Clasificación de esfuerzos}

El \textit{esfuerzo normal} es aquel cuya dirección es perpendicular al plano de la
sección transversal sobre la que actúa. Puede ser de tracción, cuando tiende a alargar
el cuerpo, o de compresión, cuando tiende a acortarlo. En este experimento, la fuerza
aplicada sobre el resorte y la liga actúa en dirección axial, por lo que el esfuerzo
resultante es de naturaleza normal \cite{halliday}.

El \textit{esfuerzo tangencial} (o cortante) actúa en dirección paralela al plano de la
sección transversal y tiende a producir un deslizamiento entre planos adyacentes del
material. El \textit{esfuerzo de torsión} es un caso particular del esfuerzo tangencial
que surge cuando se aplica un momento torsor sobre el eje longitudinal del cuerpo;
su distribución sobre la sección transversal no es uniforme, sino proporcional a la
distancia al eje de giro. Aunque ambos son de naturaleza tangencial, el esfuerzo
cortante puro se asocia a fuerzas transversales, mientras que el de torsión se origina
exclusivamente por momentos de giro \cite{serway}.

\subsection{Esfuerzo de fluencia, esfuerzo límite y módulo de elasticidad}

El \textit{módulo de elasticidad en tracción o compresión} es precisamente el módulo de
Young $Y$ definido en la Ec.~\eqref{eq:hooke}: representa la pendiente de la región
lineal del diagrama $\sigma$--$\varepsilon$ y cuantifica la rigidez intrínseca del
material frente a deformaciones axiales \cite{halliday}.

El \textit{esfuerzo límite} $\sigma_{\text{límite}}$ (o límite elástico) es el valor
máximo del esfuerzo para el cual el material aún se comporta elásticamente, es decir,
recupera su forma original al retirar la carga. Por encima de este valor la deformación
deja de ser proporcional al esfuerzo y comienza a ser permanente.

El \textit{esfuerzo de fluencia} $\sigma_y$ es el esfuerzo al que el material comienza
a deformarse plásticamente de manera apreciable y sostenida sin necesidad de incrementar
la carga. En materiales dúctiles se identifica en el diagrama $\sigma$--$\varepsilon$
como el punto donde la curva presenta una meseta horizontal o una caída brusca tras el
límite elástico, tal como se aprecia en la Fig.~\ref{fig:diagrama_sigma_epsilon}. En la
práctica, $\sigma_y$ suele determinarse mediante el criterio del $0{,}2\%$ de
deformación permanente cuando no existe una meseta claramente definida \cite{serway}.